\documentclass[11pt]{article}

\usepackage[T1]{fontenc}
\usepackage[utf8]{inputenc}
\usepackage{lmodern}
\usepackage[a4paper,margin=27mm]{geometry}
\usepackage{amsmath,amssymb}
\usepackage{booktabs}
\usepackage{array}
\usepackage{tabularx}
\usepackage{xcolor}
\usepackage[hidelinks]{hyperref}

\hypersetup{
  pdftitle={Prevalence-Calibrated Weighting for Selected Genetic Cohorts},
  pdfauthor={i3pw methodological review}
}
\setlength{\parindent}{0pt}
\setlength{\parskip}{0.55em}
\setlength{\emergencystretch}{2em}
\renewcommand{\arraystretch}{1.15}

\newcommand{\E}{\mathbb{E}}
\newcommand{\Prb}{\mathbb{P}}
\newcommand{\Var}{\operatorname{Var}}
\newcommand{\Cov}{\operatorname{Cov}}
\newcommand{\logit}{\operatorname{logit}}
\newcommand{\expit}{\operatorname{expit}}
\newcommand{\R}{\mathbb{R}}
\newcommand{\code}[1]{\texttt{#1}}
\newcommand{\pkg}{\texttt{i3pw}}

\title{\textbf{Prevalence-Calibrated Weighting for Selected Genetic Cohorts}\\
\large The idea, theory, implementation, and evidential status of \pkg}
\author{Methodological review for statistical geneticists}
\date{11 August 2026}

\begin{document}
\maketitle

\begin{abstract}
Volunteer biobanks and ascertained genetic cohorts need not have the same
distribution as their source populations.  The \pkg\ package combines weights from
a conventional participation model with an exponential tilt chosen to reproduce
externally known disease prevalences.  In the one-disease case this reduces to the
familiar case--control weights \(K/P\) for cases and
\((1-K)/(1-P)\) for controls.  In the general case it is an entropy-calibration
estimator: useful, computationally modest, and scientifically interpretable, but not
a machine for reconstructing unobserved inclusion probabilities from prevalence
alone.  This report derives the estimator, states the identifying condition, explains
its connection to ascertainment correction and label shift, reviews the software, and
separates algebraic calibration from empirical evidence.  It is written for readers
comfortable with likelihoods, generalized linear models, and GWAS, but not
necessarily with survey sampling or epidemiologic terminology.
\end{abstract}

\textbf{Review scope.}
This report describes the local version-0.3.0 worktree reviewed on the date above.
It evaluates the scientific contract as well as the code structure.  Numerical
results reported below were regenerated after the simulator corrections and frozen
in \code{report/validation\_results.tsv}; the exact local runtime is recorded in
\code{report/validation\_environment.txt}.  They remain fixed-seed synthetic evidence
rather than independent validation.

\section{Executive assessment}

The central idea is sound and classical in its mathematical form.  A register margin
is used as a moment constraint; among all participant weights satisfying that
constraint, \pkg\ chooses the weights closest in Kullback--Leibler divergence to a
user-supplied base weighting.  This joins survey calibration
\cite{deville1992,kott2010}, entropy balancing \cite{hainmueller2012}, density-ratio
models \cite{qin1998}, and the case--control correction familiar in statistical
genetics \cite{manski1977,lee2011}.

Its scientific claim must nevertheless be kept narrow.  Exact agreement with an
input prevalence is an identity, not validation.  The weights equal full-population
inverse-probability weights only under a stated density-ratio model.  If selection
also changes ancestry, severity, genotype, or exposure distributions within disease
classes, a marginal prevalence cannot repair those joint distributions.

\begin{table}[htbp]
\centering
\caption{The shortest accurate description of \pkg.}
\label{tab:executive}
\small
\begin{tabularx}{\textwidth}{>{\raggedright\arraybackslash}p{0.29\textwidth}X}
\toprule
Question & Answer \\
\midrule
What is computed? &
Normalized minimum-divergence weights matching specified population moments. \\
What information is added? &
External prevalences or other population margins, optionally layered on weights from
a participation model. \\
What is identified? &
The target population law if its density ratio relative to participants is the base
ratio times the chosen exponential tilt; otherwise an information projection, not
the true inclusion probabilities. \\
What is \(\lambda\)? &
The fitted coefficient in the log population-to-participant density ratio.  It is
not, in general, the coefficient of outcome in a participation logit. \\
Which IP weight targets the full population? &
\(1/\pi\), where \(\pi=\Prb(S=1\mid X,Y)\).  The inverse odds
\((1-\pi)/\pi\) target nonparticipants. \\
Does matching prevalence validate the method? &
No.  A held-out margin can diagnose failure, but its standardized mean difference is
a heuristic discrepancy, not a formal Sargan--Hansen \(J\)-test. \\
How strong is the present evidence? &
Theoretical links, unit tests, and synthetic studies.  The package currently has no
reported external-cohort validation; its evidential base is simulation-only. \\
\bottomrule
\end{tabularx}
\end{table}

\begin{figure}[htbp]
\centering
\fbox{\begin{minipage}{0.91\textwidth}
\centering
\begin{tabularx}{0.95\linewidth}{
>{\centering\arraybackslash}X c >{\centering\arraybackslash}X}
sampling-frame model & & register moments \(t\) \\
\(\Prb(S=1\mid X)\to d(X)\) & \(\searrow\quad\swarrow\) &
\(\E_T\{g(Y)\}\) \\
\multicolumn{3}{c}{
\(\boxed{w_i\propto d_i\exp\{\lambda^\mathsf{T}(g_i-t)\}}\)} \\
\multicolumn{3}{c}{\(\downarrow\)} \\
\multicolumn{3}{c}{weighted participant outcomes and target estimands}
\end{tabularx}

\vspace{0.7em}
Population margins deliberately withheld from the solve
\(\longrightarrow\) balance diagnostics that can fail.
\end{minipage}}
\caption{The two information channels in prevalence-calibrated weighting.  The
participation model supplies individual variation visible in frame covariates; the
register anchors selected population moments.}
\label{fig:pipeline}
\end{figure}

\section{The problem in a geneticist's language}

Let \(S=1\) denote participation, \(X\) variables available on the sampling frame,
and \(Y\) one or more diagnoses observed in participants.  Let \(T\) denote the
source population to which inference is intended to apply.  A volunteer cohort
usually samples from
\(\Prb(X,Y\mid S=1)\), whereas the scientific target is a functional of
\(\Prb_T(X,Y)\).  This is a distribution-shift problem.

For a geneticist, the simplest instance is ordinary case--control ascertainment.
If cases constitute \(P\) of the analysed sample but \(K\) of the population, a
population mean is recovered by giving cases relative weight \(K/P\) and controls
relative weight \((1-K)/(1-P)\).  \pkg\ keeps that calculation but permits:
\begin{enumerate}
\item a nonuniform starting weight \(d(X)\);
\item several prevalence, comorbidity, or stratum-specific constraints; and
\item downstream weighted means and liability-scale calculations.
\end{enumerate}
The price of this generality is that one must say exactly what selection model these
constraints identify.

\subsection{A small epidemiologic dictionary}

\begin{table}[htbp]
\centering
\caption{Terms used in this report and their statistical-genetic analogues.}
\label{tab:dictionary}
\small
\begin{tabularx}{\textwidth}{>{\raggedright\arraybackslash}p{0.24\textwidth}
>{\raggedright\arraybackslash}p{0.32\textwidth}X}
\toprule
Term & Meaning here & Familiar analogue \\
\midrule
Target population &
The people about whom the estimand is defined, including nonparticipants. &
The population whose disease prevalence \(K\) enters a liability transform. \\
Participation probability &
\(\pi(x,y)=\Prb(S=1\mid X=x,Y=y)\). &
An ascertainment or recruitment probability. \\
Positivity &
Every target-relevant type has nonzero chance to appear. &
Both cases and controls, and every required stratum, must be represented. \\
Calibration &
Choose weights so specified weighted sample moments equal known population moments. &
Force the weighted case fraction to equal \(K\). \\
Density ratio &
\(d\Prb_T/d\Prb(\cdot\mid S=1)\), the weight that transports one law to another. &
A sample-to-population prior or ascertainment correction. \\
Collider bias &
An association induced or distorted by conditioning on a common effect such as
participation. &
Genotype and phenotype both influence volunteering, so analysing volunteers can
associate them spuriously \cite{munafo2018,schoeler2023}. \\
\bottomrule
\end{tabularx}
\end{table}

\section{Probability setup: what full-population IPW means}

Write \(Z=(X,Y)\) and
\begin{equation}
\pi(z)=\Prb(S=1\mid Z=z), \qquad p_S=\Prb(S=1).
\label{eq:pi}
\end{equation}
Bayes' rule gives the population-to-participant density ratio
\begin{equation}
r_T(z)
=\frac{d\Prb_T}{d\Prb_1}(z)
=\frac{p_S}{\pi(z)},
\qquad
\Prb_1(\,\cdot\,)=\Prb_T(\,\cdot\mid S=1).
\label{eq:true-ratio}
\end{equation}
Consequently, for any integrable target variable \(h(Z)\),
\begin{equation}
\mu_T=\E_T\{h(Z)\}
=\frac{\E_1\{h(Z)/\pi(Z)\}}{\E_1\{1/\pi(Z)\}}.
\label{eq:hajek}
\end{equation}
Equation~\eqref{eq:hajek} is the self-normalized, or H\'ajek, form of
Horvitz--Thompson weighting \cite{horvitz1952,hajek1971}.  The unknown sampling fraction
\(p_S\) cancels.  It does not follow that \(\pi(z)\) can be learned from participant
data and one marginal prevalence: that is the identification problem addressed
below.

\subsection{Inverse probability is not inverse odds}

There are two legitimate transports, but they have different destinations.  From
Bayes' rule,
\begin{align}
\frac{d\Prb_T}{d\Prb_1}(z)
&\propto \frac{1}{\pi(z)}, \label{eq:full-ipw}\\
\frac{d\Prb_0}{d\Prb_1}(z)
&\propto \frac{1-\pi(z)}{\pi(z)},
\qquad
\Prb_0(\,\cdot\,)=\Prb_T(\,\cdot\mid S=0).
\label{eq:odds-transport}
\end{align}
Thus \(1/\pi\) targets the full source population, whereas inverse odds target
nonparticipants.  Logistic modelling does not make the two targets identical.
When participation is rare, \(1-\pi(z)\approx1\), so the normalized weights can
be close; this is an approximation, not an algebraic equivalence.

\section{The one-outcome calculation}

Let \(Y\in\{0,1\}\), let
\begin{equation}
P=\Prb(Y=1\mid S=1),\qquad K=\Prb_T(Y=1),
\label{eq:sample-target-prevalence}
\end{equation}
and begin from uniform base weights.  An exponential tilt gives cases relative
weight \(\exp(\lambda)\) and controls weight one.  Requiring the weighted case
fraction to equal \(K\) yields
\begin{equation}
\frac{P\exp(\lambda)}
{P\exp(\lambda)+(1-P)}=K
\quad\Longrightarrow\quad
\lambda=\logit(K)-\logit(P).
\label{eq:binary-lambda}
\end{equation}
After an irrelevant common normalization, the class weights are
\begin{equation}
d_1=\frac{K}{P},
\qquad
d_0=\frac{1-K}{1-P}.
\label{eq:class-weights}
\end{equation}
For \(P=0.20\) and \(K=0.40\), equation~\eqref{eq:binary-lambda} gives
\(\lambda=0.9808\); the relative weights are \(2.00\) for a case and \(0.75\)
for a control.  This is worth checking by hand before trusting a larger solve.
This is the classical choice-based-sampling result \cite{manski1977} and the
core ascertainment factor appearing in liability-scale corrections
\cite{lee2011}.

The interpretation of \(\lambda\) matters.  Let
\(\pi_y=\Prb(S=1\mid Y=y)\).  Equations~\eqref{eq:true-ratio} and
\eqref{eq:binary-lambda} imply
\begin{equation}
\lambda=\log\!\left(\frac{\pi_0}{\pi_1}\right).
\label{eq:lambda-density}
\end{equation}
It is therefore a coefficient in the \emph{population-to-participant density
ratio}.  Suppose instead that participation obeys
\(\logit(\pi_y)=\alpha+\gamma y\).  Then
\begin{equation}
\lambda
=-\gamma+
\log\!\left\{\frac{1+\exp(\alpha+\gamma)}
{1+\exp(\alpha)}\right\}
=-\gamma+\log\!\left(\frac{1-\pi_0}{1-\pi_1}\right).
\label{eq:gamma-relation}
\end{equation}
Only under rare participation is the second term negligible, giving
\(\lambda\approx-\gamma\).  Calling \(\lambda\) the participation-logit
coefficient would therefore reverse a sign and omit a non-negligible term in
many cohort settings.

\section{General entropy calibration}

Let \(g_i=g(Y_i,X_i)\in\R^k\) be calibration features, with externally known target
\begin{equation}
t=\E_T\{g(Y,X)\}.
\label{eq:target-moment}
\end{equation}
Examples include diagnosis indicators, pairwise comorbidity indicators, stratum
shares, and diagnosis-by-stratum products.  Let \(d_i\geq0\) be normalized base
weights.  Uniform \(d_i\) give pure prevalence calibration; estimated
\(1/\widehat\pi(X_i)\) supply a covariate correction first.

\subsection{Primal and dual problems}

For exact calibration, \pkg\ solves the empirical information projection
\begin{align}
\min_{w_1,\ldots,w_n}\quad&
\sum_{i=1}^n w_i\log\!\left(\frac{w_i}{d_i}\right)
\label{eq:primal}\\
\text{subject to}\quad&
w_i\geq0,\qquad \sum_{i=1}^n w_i=1,\qquad
\sum_{i=1}^n w_i g_i=t. \notag
\end{align}
This is a Kullback--Leibler calibration estimator
\cite{csiszar1975,deville1992,hainmueller2012}.  Its dual is
\begin{equation}
\widehat\lambda
=\arg\min_{\lambda\in\R^k}
\log\!\left[
\sum_{i:d_i>0}d_i
\exp\{\lambda^\mathsf{T}(g_i-t)\}
\right],
\label{eq:dual}
\end{equation}
and the resulting weights are
\begin{equation}
w_i(\lambda)=
\frac{d_i\exp\{\lambda^\mathsf{T}(g_i-t)\}}
{\sum_{j:d_j>0}d_j\exp\{\lambda^\mathsf{T}(g_j-t)\}}.
\label{eq:tilted-weights}
\end{equation}
Rows with \(d_i=0\) remain outside the support and receive zero weight.  The
optimization is only \(k\)-dimensional even when \(n\) is large.

Empirical likelihood and entropy balancing are close relatives, not the same
optimization.  Both impose moment restrictions and belong to generalized empirical
likelihood families.  Standard empirical likelihood maximizes
\(\sum_i\log w_i\) and produces reciprocal-affine weights; the criterion in
\eqref{eq:primal} produces exponential weights.  Results from one should not be
attributed mechanically to the other \cite{qinlawless1994,qin1998}.

\subsection{Ridge stabilization}

With ridge parameter \(\rho>0\), the implemented dual becomes
\begin{equation}
\widehat\lambda_\rho
=\arg\min_\lambda
\left[
\log\sum_{i:d_i>0}d_i
\exp\{\lambda^\mathsf{T}(g_i-t)\}
+\frac{\rho}{2}\|\lambda\|_2^2
\right].
\label{eq:ridge-dual}
\end{equation}
Its gradient and Hessian are
\begin{align}
\nabla\ell_\rho(\lambda)
&=\E_{w(\lambda)}(g)-t+\rho\lambda,
\label{eq:ridge-gradient}\\
\nabla^2\ell_\rho(\lambda)
&=\Cov_{w(\lambda)}(g)+\rho I_k.
\label{eq:ridge-hessian}
\end{align}
At the solution,
\begin{equation}
\E_{w(\widehat\lambda_\rho)}(g)-t
=-\rho\widehat\lambda_\rho.
\label{eq:ridge-residual}
\end{equation}
Ridge calibration is therefore deliberately approximate.  It trades bias in the
specified margins against less concentrated weights; it must not be reported as exact
calibration.

\subsection{Feasibility and positivity}

For a finite exact solution, \(t\) must lie in the relative interior of the convex
hull of the observed \(g_i\) having positive base weight.  A binary target strictly
between zero and one is unreachable if the sample contains only cases or only
controls.  Near a hull boundary, \(\|\widehat\lambda\|\) and the largest weight can
explode even when the optimizer reports success.  Redundant constraints can also make
\(\lambda\) nonunique.  These are data-support problems, not numerical inconveniences
that a different optimizer can wish away.

\section{Identification: when calibration is population weighting}

The empirical solve always defines weights if its constraints are feasible.  The
scientific question is when those weights recover the target law.

\paragraph{Density-ratio assumption.}
Suppose the true population-to-participant ratio can be represented as
\begin{equation}
r_T(X,Y)
=c\,d(X)\exp\{\theta^\mathsf{T}g(X,Y)\},
\label{eq:identification-model}
\end{equation}
where \(c\) normalizes the ratio.  Here \(d(X)\) is the population-to-participant
correction already supplied by the base model, and \(g\) spans the remaining
selection pattern.  If the target moments are correct, positivity holds, and the
moment map is identifiable, then the population moment equation
\begin{equation}
\frac{\E_1[d(X)\exp\{\theta^\mathsf{T}g\}g]}
{\E_1[d(X)\exp\{\theta^\mathsf{T}g\}]}
=t
\label{eq:population-moment-equation}
\end{equation}
identifies \(\theta\); the sample solution estimates it.  Under
\eqref{eq:identification-model}, the calibrated weights are proportional to the true
full-population IP weights.

This statement is the right meaning of \emph{prevalence informs the weights}.  It does
not say that prevalence nonparametrically identifies \(\pi(X,Y)\).  Rather, prevalence
identifies the low-dimensional coefficient of a density-ratio family chosen in
advance.

\subsection{What happens under misspecification}

If equation~\eqref{eq:identification-model} is false, the solve still returns the
member of the exponential family that matches \(t\) and is closest to the base
distribution in KL divergence.  Let \(r_\lambda\) denote that fitted ratio.  For a
downstream variable \(h\), its transport bias is
\begin{equation}
\E_1\!\left[\{r_\lambda(Z)-r_T(Z)\}h(Z)\right].
\label{eq:transport-bias}
\end{equation}
Calibration gives no general reason for equation~\eqref{eq:transport-bias} to vanish.
It does vanish for functions in the affine span of the constrained moments:
\begin{equation}
h(Z)=a+b^\mathsf{T}g(Z)
\quad\Longrightarrow\quad
\E_{w_\lambda}(h)=a+b^\mathsf{T}t.
\label{eq:span-property}
\end{equation}
This explains both the method's strength and its limit.  Anchored quantities are
right by construction; transfer to unanchored quantities depends on how well the
chosen \(g\) captures the relevant distribution shift.

\subsection{Label shift and case mix}

With uniform base weights, a categorical \(Y\), and \(g(Y)\) containing a full set
of class indicators, the assumption reduces to
\begin{equation}
\Prb_T(X\mid Y)=\Prb_1(X\mid Y).
\label{eq:label-shift}
\end{equation}
That is, only class proportions change.  This is prior-probability or label shift
\cite{saerens2002,lipton2018}.  It is plausible for deliberately sampled
case--control data when recruitment is random within case status.  It is doubtful
when, for example, young severe cases enter at a different rate from old mild cases.

For \(Q\) comorbid binary diagnoses, matching only the \(Q\) marginal prevalences is
not full label-shift correction over the \(2^Q\) joint disease patterns.  If
participation treats a comorbid case differently from the product of its marginal
diagnoses, the corresponding joint prevalence or interaction must be constrained.

A known \(K=\Prb_T(Y=1)\) sets the number of cases, not their composition:
\(\Prb_T(X\mid Y=1)\) remains unidentified.  If population prevalences are available
within sex, birth cohort, ancestry, region, or severity strata, interactions such as
\begin{equation}
g_{aq}(X,Y)=\mathbf{1}(A=a)Y_q
\label{eq:stratified-feature}
\end{equation}
can represent that structure.  The number of constraints then grows quickly, so
richer identification consumes positivity and effective sample size.  There is no
free lunch; the convex hull sends an invoice.

\section{Target estimands in statistical genetics}

Weighting changes a probability law.  Whether that repairs a scientific analysis
depends on whether the analysis is a functional of that law and whether the fitted
density ratio is correct.

\begin{table}[htbp]
\centering
\caption{What prevalence calibration can and cannot establish for common genetic
analyses.}
\label{tab:estimands}
\footnotesize
\begin{tabularx}{\textwidth}{
>{\raggedright\arraybackslash}p{0.22\textwidth}
>{\raggedright\arraybackslash}p{0.33\textwidth}X}
\toprule
Estimand & When the weighting helps & Principal caveat \\
\midrule
Disease prevalence &
An exact anchor reports the supplied population prevalence. &
This is not an empirical result; uncertainty in the external prevalence remains. \\
Population mean or absolute risk &
Consistent if the fitted ratio transports the participant distribution of the
measured variable to the target population. &
Marginal disease calibration alone does not fix within-disease selection on severity
or prognosis. \\
Allele frequency or biomarker mean &
Can transfer when its relation to the constrained features is stable, or when the
full density-ratio model is correct. &
A known disease prevalence contains no direct information about genotype or biomarker
composition within disease classes. \\
Logistic GWAS coefficient &
Under selection depending on case status alone, the odds-ratio slope is invariant to
case--control sampling \cite{prentice1979}. &
When genotype or its causes and phenotype both affect participation, conditioning on
\(S=1\) can create collider bias; a marginal \(K\) does not identify the required
joint correction \cite{munafo2018,schoeler2023}. \\
Liability-scale \(R^2\) or heritability &
Case/control weights recover the population case fraction before the
observed-to-liability transform; this parallels classical ascertainment correction
\cite{dempster1950,lee2011}. &
The result inherits the liability-threshold model and the assumed selection law.
The current pairwise implementation is quadratic in sample size. \\
Causal effect &
Weights can address selection into the analysed sample. &
They do not by themselves address exposure confounding, measurement error, treatment
assignment, or interference. \\
\bottomrule
\end{tabularx}
\end{table}

\subsection{Participation as a collider}

Suppose genotype \(G\) or an exposure \(E\) influences participation and disease \(Y\)
also influences participation.  Conditioning on volunteers opens a path between the
causes of \(S\), potentially distorting \(G\)--\(Y\) or \(E\)--\(Y\) associations.
This is not repaired merely by restoring \(\Prb(Y=1)=K\), because infinitely many
joint laws \(\Prb_T(G,Y)\) share the same marginal \(K\).  Volunteer-weighting studies
in UK Biobank illustrate both the importance and the difficulty of modelling this
selection \cite{schoeler2023,vanalten2024,vanalten2025}.

The distinction from Prentice--Pyke invariance is useful:
\begin{itemize}
\item if \(S\) depends only on \(Y\), a correctly specified logistic disease slope
can remain consistent without prevalence weighting; and
\item if \(S\) depends jointly on \(G\), \(X\), and \(Y\), population associations
require a correction rich enough for that joint selection.
\end{itemize}
Prevalence calibration can be one component of that correction, but it cannot replace
the other components.

\subsection{Liability-scale quantities}

For population prevalence \(K\), threshold
\(t_K=\Phi^{-1}(1-K)\), and \(z_K=\phi(t_K)\), a common
observed-to-liability conversion is
\begin{equation}
R_L^2
=R_O^2\frac{K(1-K)}{z_K^2}.
\label{eq:obs-liability}
\end{equation}
If \(R_O^2\) was estimated in a case--control sample with case fraction \(P\), the
Lee et al.\ ascertainment factor gives
\begin{equation}
R_L^2
=R_{O,\mathrm{asc}}^2
\times\frac{K(1-K)}{z_K^2}
\times\frac{K(1-K)}{P(1-P)}.
\label{eq:lee-transform}
\end{equation}
Alternatively, one may first use equation~\eqref{eq:class-weights} to estimate the
population-standardized observed-scale moment and then apply
equation~\eqref{eq:obs-liability}.  Their agreement is expected under pure
case/control ascertainment.  Neither route protects against recruitment depending on
latent liability within case status unless that dependence is modelled.

\section{Uncertainty: what is random?}

There are at least four uncertainty sources: sampling participants, estimating the
base participation model, fitting the tilt, and estimating or transporting the
external targets.  A standard error that conditions on all weights and targets treats
only the first.

\subsection{Fixed-weight and calibration-aware linearization}

For normalized weights \(\widetilde w_i=w_i/\sum_jw_j\), the fixed-weight variance of
\(\widehat\mu=\sum_i\widetilde w_i h_i\) is estimated by
\begin{equation}
\widehat{\Var}_{\mathrm{fixed}}(\widehat\mu)
=\sum_{i=1}^n\widetilde w_i^2(h_i-\widehat\mu)^2.
\label{eq:fixed-var}
\end{equation}
This ignores that the tilt was re-estimated to meet the constraints.  The
calibration-aware linearization implemented in \pkg\ instead uses
\begin{align}
\widehat\beta_\rho
&=\{\widehat{\Cov}_w(g)+\rho I\}^{-1}
\widehat{\Cov}_w(g,h),
\label{eq:if-beta}\\
e_i
&=h_i-\widehat\mu
-\widehat\beta_\rho^\mathsf{T}(g_i-\bar g_w),
\label{eq:if-residual}\\
\widehat{\Var}_{\mathrm{cal}}(\widehat\mu)
&=\sum_{i=1}^n\widetilde w_i^2e_i^2.
\label{eq:cal-var}
\end{align}
The software computes equation~\eqref{eq:if-beta} by penalized least squares rather
than explicitly inverting a possibly ill-conditioned covariance matrix.

If \(\rho=0\) and \(h\) is an exactly constrained column, the residual in
equation~\eqref{eq:if-residual} is zero up to numerical tolerance.  Its conditional
sampling standard error is consequently zero: the estimator always prints the fixed
target.  For \(\rho>0\), equation~\eqref{eq:ridge-residual} shows that the target is not
fixed, \(\widehat\beta_\rho\) is shrunk, and the anchored margin retains uncertainty.
Using the exact-calibration projection after a ridge solve would be wrong.

\subsection{External targets, fitted bases, and dependence}

Zero conditional standard error is not zero scientific uncertainty.  If a registry
target \(\widehat t\) has covariance \(V_t\), a first-order combination for an
independent external source is
\begin{equation}
\Var(\widehat\mu)
\approx \Var_{\mathrm{participants}}(\widehat\mu\mid\widehat t)
+D_tV_tD_t^\mathsf{T},
\qquad
D_t=\frac{\partial\mu}{\partial t^\mathsf{T}}.
\label{eq:target-uncertainty}
\end{equation}
Temporal mismatch, outcome-definition mismatch, ancestry mismatch, and overlap
between registry and cohort create additional bias or covariance that
equation~\eqref{eq:target-uncertainty} does not solve.

The package also offers resampling for its simulation wrapper.  A scientifically
complete bootstrap should re-fit every estimated stage that would change under
repeated sampling, including the base model when applicable.  Rare anchors can make
bootstrap replicates infeasible; silently discarding precisely those replicates
truncates the distribution and narrows intervals.  Family or household relatedness
requires cluster-level resampling or a cluster-robust influence calculation.  The
current simple mean interface does not automate those genetic-cohort complications.

\subsection{Augmented IPW}

For a participant-only trait \(V\), covariates \(X\) available across the population
frame, and regression \(m(X)=\E(V\mid X)\), the package implements
\begin{equation}
\widehat\mu_{\mathrm{AIPW}}
=\frac{1}{N}\sum_{i=1}^{N}\widehat m(X_i)
+\sum_{i:S_i=1}\widetilde w_i
\{V_i-\widehat m(X_i)\}.
\label{eq:aipw}
\end{equation}
In the standard missing-data or data-integration model, this estimator can be
consistent when either the outcome regression or the participation weighting is
correct \cite{robins1994,bang2005}; flexible regressions call for sample splitting or
cross-fitting \cite{chernozhukov2018}.  This is not a blanket double-robustness theorem
for prevalence
calibration.  Similarly, the result that entropy balancing is doubly robust in
particular causal models \cite{zhao2017} does not prove that arbitrary
outcome-dependent volunteer selection is repaired by a few prevalence margins.

\section{Diagnostics and falsifiability}

\subsection{Quantities that diagnose the solve}

\begin{table}[htbp]
\centering
\caption{Diagnostics, their interpretation, and their limits.}
\label{tab:diagnostics}
\footnotesize
\begin{tabularx}{\textwidth}{
>{\raggedright\arraybackslash}p{0.23\textwidth}
>{\raggedright\arraybackslash}p{0.34\textwidth}X}
\toprule
Diagnostic & What it reveals & What it cannot reveal \\
\midrule
Optimizer convergence &
Whether the numerical dual solve terminated satisfactorily. &
Whether the density-ratio family is scientifically correct. \\
Calibration residual &
Whether exact targets were reached, or how far ridge calibration missed them. &
Validity: an exact residual is expected by construction. \\
Kish effective sample size \cite{kish1965} &
\(\mathrm{ESS}=(\sum_iw_i)^2/\sum_iw_i^2\), a concise measure of concentration. &
A universal variance or information measure for every downstream analysis. \\
Weight range and top-1\% mass &
Dependence on a small number of observations. &
Bias from omitted selection variables. \\
Support counts &
Whether each required class or stratum is represented. &
Adequate overlap within unmeasured subgroups. \\
Held-out SMD &
Whether a known margin excluded from the solve is reproduced after weighting. &
A formal hypothesis test, or proof that all unobserved margins are balanced. \\
\bottomrule
\end{tabularx}
\end{table}

For a held-out variable \(a\) with known population mean \(t_a\), \pkg\ reports
\begin{equation}
\mathrm{SMD}_a
=\frac{\sum_i\widetilde w_i a_i-t_a}{s_{a,\mathrm{sample}}}.
\label{eq:smd}
\end{equation}
The often-used rule \(|\mathrm{SMD}|<0.1\) is a descriptive convention
\cite{austin2009}.  A held-out failure is valuable: it refutes the proposition that
the fitted weighting transports that margin.  A small SMD is only partial positive
evidence.

Calling equation~\eqref{eq:smd} a Sargan--Hansen \(J\)-test would overstate it
\cite{sargan1958,hansen1982}.
A formal overidentification test needs a joint covariance estimate, a quadratic test
statistic, degrees of freedom, and an asymptotic reference distribution while
accounting for estimated parameters.  The present diagnostic borrows the
\emph{logic} of unused restrictions, not that formal test.

\subsection{An honest validation design}

Known margins should be divided before fitting:
\begin{enumerate}
\item use a prespecified subset to calibrate;
\item keep other registry quantities genuinely held out;
\item fit the tilt on one wave or fold and use the stored tilt on a new wave when
possible; and
\item evaluate downstream quantities not algebraically determined by the anchors.
\end{enumerate}
Selecting constraints after inspecting held-out failures converts validation into
training.  Reporting only anchored prevalence error produces impressive zeros and no
evidence whatsoever.

\section{Software review}

\subsection{Architecture and public contract}

The reviewed source declares version 0.3.0 and an alpha development status.  Its
design has a sensible separation between a real-cohort interface and simulation
machinery.

\begin{table}[htbp]
\centering
\caption{Main software components in the reviewed \pkg\ worktree.}
\label{tab:architecture}
\footnotesize
\begin{tabularx}{\textwidth}{
>{\raggedright\arraybackslash}p{0.24\textwidth}
>{\raggedright\arraybackslash}p{0.28\textwidth}X}
\toprule
Component & Public role & Review \\
\midrule
\code{calibrate} and \code{CalibrationFit} &
Real-cohort entry point using participant outcomes, external targets, optional base
weights, strata, joint moments, and held-out margins. &
The correct conceptual front door; it avoids requiring simulated ground truth. \\
\code{entropy\_balance} and \code{apply\_tilt} &
Core dual solve and transfer of a fitted tilt to new rows. &
Small convex problem, log-scale stabilization, explicit diagnostics, and preservation
of zero-weight support. \\
\code{lasso\_propensity} and \code{lasso\_ipw} &
Covariate-only participation baseline. &
Useful comparator, but a real application needs frame-level participation labels and
compatible covariates for participants and nonparticipants. \\
\code{calibration\_ipw} and \code{monte\_carlo} &
Simulation wrappers with known population truth. &
Appropriate for method development, not the interface to advertise as real-data
inference. \\
\code{balance} and \code{uncertainty} &
Held-out discrepancies, effective sample size, linearization, bootstrap, and target
sensitivity. &
Good instincts; target uncertainty, clustering, and a formal specification test
remain user responsibilities. \\
\code{aipw} &
Outcome-regression augmentation for a participant-only trait when \(X\) is known for
the population frame. &
Cross-fitting is available; robustness claims still require the usual missing-data
conditions. \\
\code{liability} &
Threshold-model simulations and weighted pairwise moment estimators. &
Scientifically relevant, but the explicit \(n\times n\) similarity matrix does not
scale to biobank sample sizes. \\
\bottomrule
\end{tabularx}
\end{table}

\subsection{Scientific contracts that must remain explicit}

\begin{table}[htbp]
\centering
\caption{High-impact semantic and numerical contracts in the corrected design.}
\label{tab:contracts}
\footnotesize
\begin{tabularx}{\textwidth}{
>{\raggedright\arraybackslash}p{0.25\textwidth}
>{\raggedright\arraybackslash}p{0.33\textwidth}X}
\toprule
Issue & Correct contract & Consequence \\
\midrule
Meaning of the tilt &
\(\lambda\) belongs to the log population-to-participant density ratio, as in
equations~\eqref{eq:lambda-density}--\eqref{eq:gamma-relation}. &
Do not report it as the outcome coefficient in a participation logit. \\
Target of base weights &
\(1/\widehat\pi\) transports participants to the full population;
\((1-\widehat\pi)/\widehat\pi\) transports them to nonparticipants. &
The inverse scheme is the full-population default; an odds option needs an explicit
different target. \\
Simulation oracle &
A full-population oracle is the complete simulated test-fold mean, unavailable in
real data. &
An inverse-odds mixture is not a full-population oracle and should not be labelled one. \\
Ridge calibration &
The achieved residual equals \(-\rho\widehat\lambda_\rho\), and the influence
projection contains the same \(\rho\). &
Only exact \(\rho=0\) anchors have zero conditional standard error. \\
Zero base weights &
Zero-weight rows are excluded from the exponential shift and remain zero. &
An extreme feature value outside the positive support must not destabilize all
weights. \\
Held-out balance &
The SMD is a heuristic discrepancy using an unused margin. &
It may be called a falsification diagnostic, not a formal \(J\)-test. \\
Synthetic selection &
The generator uses Bernoulli logistic participation with coefficients chosen so
expected sample count and outcome totals meet configured targets. &
Realized sample size and prevalences are random; the mechanism is not fixed-size
probability-proportional-to-size sampling. \\
\bottomrule
\end{tabularx}
\end{table}

\subsection{Computational properties}

Let \(n\) be the number of participants with positive base weight, \(k\) the number
of constraints, \(p\) the number of frame covariates, \(I\) optimizer iterations,
and \(B\) bootstrap replicates.

\begin{table}[htbp]
\centering
\caption{Leading computational costs.  These are algorithmic orders, not timing
benchmarks.}
\label{tab:complexity}
\footnotesize
\begin{tabularx}{\textwidth}{
>{\raggedright\arraybackslash}p{0.25\textwidth}
>{\raggedright\arraybackslash}p{0.22\textwidth}
>{\raggedright\arraybackslash}p{0.18\textwidth}X}
\toprule
Operation & Approximate time & Approximate memory & Comment \\
\midrule
Entropy dual &
\(O(Ink)\) &
\(O(nk)\) for features &
L-BFGS avoids forming a dense \(k\times k\) Hessian. \\
Stratified design &
\(O(nk)\) &
\(O(nk)\) &
\(k\) can be roughly strata times outcomes, plus share and interaction constraints. \\
LASSO participation model &
Solver- and CV-dependent &
\(O(np)\) before expansion &
Pairwise interactions increase the design toward \(O(np^2)\). \\
Calibration bootstrap &
\(O(BInk)\) plus optional base refits &
One replicate plus stored output &
Parallelization is natural, but rare anchors can cause infeasible replicates. \\
Liability pairwise moment &
\(O(n^2p)\) to form \(XX^\mathsf{T}\) &
\(O(n^2)\) &
This is the principal barrier to biobank-scale use.  Blocked matrix products or
summary-statistic identities are needed. \\
Synthetic data generator &
Repeated \(O(Nq)\) selection objectives, plus covariance work &
\(O(Np+Nq+p^2)\) &
Useful for controlled experiments; not normally a production bottleneck. \\
\bottomrule
\end{tabularx}
\end{table}

For the core calibration, \(k\) is normally small and the implementation is
reasonable.  The main practical failure is statistical rather than computational:
adding many sparse stratum-by-disease constraints can exhaust support before it
exhausts RAM.  The liability module is the exception; materializing \(n^2\) pairwise
similarities is untenable at modern biobank scale.

\subsection{Remaining release gaps}

\begin{table}[htbp]
\centering
\caption{Important gaps that remain after correcting the scientific semantics.}
\label{tab:release-gaps}
\footnotesize
\begin{tabularx}{\textwidth}{
>{\raggedright\arraybackslash}p{0.25\textwidth}
>{\raggedright\arraybackslash}p{0.34\textwidth}X}
\toprule
Gap & Consequence & Recommended next step \\
\midrule
Real-data resampling interface &
The bootstrap and prevalence-sensitivity helpers take the simulation
\code{Dataset}; users of \code{calibrate} must assemble equivalent refits themselves. &
Provide a callback- or array-based bootstrap that refits the base, tilt, and estimand
without simulated ground truth. \\
External-target and cluster variance &
\code{CalibrationFit.mean} treats external targets as fixed and independent units as
the sampling units. &
Accept target covariance and cluster identifiers, or document a supported
cluster-bootstrap recipe. \\
Liability scaling &
The explicit \(n\times n\) similarity matrix is incompatible with biobank-scale
\(n\). &
Use blocked accumulations, randomized trace methods, or algebraic sufficient
statistics that avoid storing pairwise matrices. \\
Rendered documentation &
The corrected distributions include the Markdown theory, examples, LaTeX report, and
citation record, but there is no rendered versioned site or link checker. &
Publish versioned HTML/PDF documentation and check external and internal links in CI. \\
Frozen environment and automated evidence &
A citation file, structured post-correction results, and a local environment manifest
are now present, and CI smoke-tests one packaged example; there is no cross-platform
dependency lock, replicated-validation job, or comparison of outputs with the frozen
table. &
Record a locked environment and add a reproducible benchmark job with explicit
tolerances and retained logs. \\
\bottomrule
\end{tabularx}
\end{table}

\section{Evidence review}

The mathematical ingredients of \pkg\ are well established.  What is new is their
particular packaging and proposed use of register prevalences to augment
participation weights.  Evidence for that applied proposal must be judged separately
from evidence for entropy calibration itself.

\begin{table}[htbp]
\centering
\caption{Current evidence and what each layer can support.}
\label{tab:evidence}
\small
\begin{tabularx}{\textwidth}{
>{\raggedright\arraybackslash}p{0.24\textwidth}
>{\raggedright\arraybackslash}p{0.32\textwidth}X}
\toprule
Evidence layer & What it supports & What it does not support \\
\midrule
Established literature &
IPW, survey calibration, information projection, density-ratio tilting, and
case--control ascertainment correction. &
That the selected \(g\) spans selection in a particular biobank. \\
Unit tests &
Array contracts, numerical convergence, known special cases, support checks, and
selected uncertainty calculations. &
External validity, realistic misspecification, or calibrated interval coverage in a
real cohort. \\
Synthetic data generator &
Controlled comparison when outcomes and Bernoulli-logistic participation are known. &
Performance under an unknown real recruitment process. \\
Example scripts and frozen local table &
Illustrations of pure outcome selection, covariate-plus-outcome selection, collider
settings, comorbidity, and liability models; selected corrected runs are recorded in
\code{report/validation\_results.tsv}. &
Prespecified, independently replicated benchmark evidence. \\
Real-cohort study &
None is included in the present evidence base. &
No empirical claim about bias reduction in iPSYCH, UK Biobank, or another cohort
follows from the repository alone. \\
\bottomrule
\end{tabularx}
\end{table}

\subsection{Fresh post-correction smoke evidence}

Table~\ref{tab:fresh-results} records the most informative results from a fresh local
run on 11 August 2026 using Python 3.11.15, NumPy 2.4.6, SciPy 1.17.1, and
scikit-learn 1.9.0 with the corrected \pkg\ 0.3.0 worktree.  Replicated scripts
use seeds \(0,\ldots,R-1\).  These are
regression and plausibility checks, not a substitute for an independently designed
validation study.

\begin{table}[htbp]
\centering
\caption{Selected post-correction synthetic results.  Values are mean
\(\pm\) standard deviation across repetitions unless stated otherwise.}
\label{tab:fresh-results}
\footnotesize
\begin{tabularx}{\textwidth}{
>{\raggedright\arraybackslash}p{0.22\textwidth}
>{\raggedright\arraybackslash}p{0.20\textwidth}
>{\raggedright\arraybackslash}p{0.24\textwidth}X}
\toprule
Experiment & Quantity & Result & Interpretation \\
\midrule
\code{honest\_benchmark}, no \(X\) channel, \(R=12\) &
Absolute percentage error on an unanchored outcome &
Naive \(80.78\pm6.66\); calibrated LASSO base \(81.61\pm6.03\);
calibrated uniform base \(82.24\pm5.96\). &
Anchoring one disease did not transfer to another. \\
\code{honest\_benchmark}, \(X\) channel strength \(1.5\), \(R=12\) &
Worst held-out \(|\mathrm{SMD}|\); absolute error in held-out \(X_0\) mean &
LASSO base \(0.209\pm0.079\), \(0.126\pm0.072\);
uniform base \(0.482\pm0.238\), \(0.317\pm0.211\). &
When a covariate channel exists, the base participation model materially helps. \\
\code{doubly\_robust\_trait}, \(R=20\) &
Mean bias; mean absolute bias for a participant-only trait &
Naive \(-0.0905,0.1008\); LASSO IPW \(-0.0252,0.0636\);
calibration \(0.0537,0.0864\); AIPW \(-0.0009,0.0584\). &
Augmentation performed best in this particular, outcome-model-friendly design. \\
\code{monte\_carlo}, \(R=20\) &
Anchored prevalence error; calibration ESS &
Calibration error \(0.00\pm0.00\); ESS \(153\pm19\). &
The zero error is the constraint, not validation; ESS quantifies its variance cost. \\
\bottomrule
\end{tabularx}
\end{table}

The results support the package's deliberately narrow messages: exact anchors are
tautologies, cross-outcome transfer is weak, and the covariate base matters when the
data-generating process actually contains a covariate selection channel.  They do
not establish bias correction in a real cohort, interval coverage, or robustness to
an unknown selection law.  For stronger evidence:
\begin{itemize}
\item error on a constrained prevalence must be identified as a match by construction;
\item at least one outcome, stratum, or population margin must remain unanchored;
\item simulations must vary both covariate-driven and outcome-driven selection;
\item misspecified tilts, target error, weak support, and within-case selection must be
included; and
\item empirical validation needs a population quantity unavailable to the fitting
procedure, preferably in a later cohort wave or linked register.
\end{itemize}

\subsection{Recommended validation matrix}

\begin{table}[htbp]
\centering
\caption{Minimum simulation and empirical validation matrix for a methods paper.}
\label{tab:validation-matrix}
\footnotesize
\begin{tabularx}{\textwidth}{
>{\raggedright\arraybackslash}p{0.23\textwidth}
>{\raggedright\arraybackslash}p{0.30\textwidth}X}
\toprule
Axis & Levels & Primary outputs \\
\midrule
Selection law &
\(X\)-only; \(Y\)-only; additive \(X+Y\); \(X\times Y\); within-case severity;
genotype-plus-outcome. &
Bias, RMSE, interval coverage, ESS, maximum weight. \\
Anchor information &
Correct pooled prevalence; stratum-specific prevalence; joint comorbidity; perturbed
targets; partial anchors. &
Anchor residual and transfer error on held-out outcomes and margins. \\
Support &
Common, rare, and extremely rare disease; sparse ancestry/sex/age cells. &
Failure probability, convergence, weight concentration, bootstrap discard rate. \\
Base model &
Correct, misspecified, regularized, cross-fitted, and absent. &
Incremental value of the base versus the prevalence tilt. \\
Comparator &
Unweighted, covariate IPW, exact calibration, ridge calibration, oracle full-population
law in simulation. &
Bias--variance trade-off on the same target estimand. \\
Real data &
Prespecified fit margins and separate held-out register margins or recruitment wave. &
External discrepancies with uncertainty, not merely anchored matches. \\
\bottomrule
\end{tabularx}
\end{table}

\section{Practical analysis protocol}

A defensible analysis can be organized in nine steps.

\begin{enumerate}
\item \textbf{Define \(T\).}  State geography, calendar period, age range, ancestry
scope, survival condition, and disease definition.  A prevalence from a different
population is not a magic number; it is a different estimand.

\item \textbf{Define the target quantity.}  A prevalence, population mean, regression
coefficient, and liability-scale variance component need different estimating
equations.  Weighting is not the estimand.

\item \textbf{Build the base correction.}  Fit participation using variables observed
compatibly in the cohort and its sampling frame.  For the full population use
\(1/\widehat\pi(X)\), not inverse odds.  Separate training and analysis folds when a
flexible model is used.

\item \textbf{Prespecify \(g\).}  Begin with scientifically justified prevalence
margins.  Add strata or interactions only when population targets and participant
support exist.  Reserve some credible margins for validation.

\item \textbf{Fit and inspect.}  Report convergence, residuals, tilt coefficients,
support counts, ESS, range, quantiles, and top-weight mass.  For ridge, report
\(\rho\) and achieved rather than nominal margins.

\item \textbf{Try to refute the weighting.}  Evaluate held-out margins and a transferred
tilt on a later fold or wave.  Do not call an SMD threshold a \(p\)-value.

\item \textbf{Estimate uncertainty at the right level.}  Use the penalized calibration
influence adjustment corresponding to the fitted \(\rho\), or re-fit the entire
pipeline in a bootstrap.  Respect family clusters and quantify uncertainty in external
targets.

\item \textbf{Stress the assumptions.}  Perturb prevalence, alter trimming or ridge,
coarsen strata, and examine plausible within-case selection.  Trimming after exact
calibration breaks the constraint and must be followed by reporting the new residual.

\item \textbf{Report weighted and unweighted analyses.}  Large changes deserve
explanation; small changes can reflect either little bias or an uninformative weight.
Neither is established by the point estimate alone.
\end{enumerate}

\begin{table}[htbp]
\centering
\caption{Items that should accompany a published \pkg\ analysis.}
\label{tab:reporting}
\small
\begin{tabularx}{\textwidth}{>{\raggedright\arraybackslash}p{0.30\textwidth}X}
\toprule
Item & Minimum disclosure \\
\midrule
Target population &
Eligibility, time, place, ancestry handling, outcome definition, and source of every
external margin. \\
Base model &
Training frame, predictors, missing-data handling, validation, clipping, and whether
weights are \(1/\pi\) or inverse odds. \\
Calibration model &
All \(g\) columns and targets, exact or ridge objective, \(\rho\), strata, joint
moments, and any post-fit trimming. \\
Support and concentration &
Cases and controls per cell, ESS, weight quantiles, maximum, and top-1\% mass. \\
Validation &
Which margins were held out before fitting, their discrepancies, and whether the tilt
was transferred to independent data. \\
Inference &
Influence or bootstrap procedure, refitted stages, clustering unit, failed
replicates, and treatment of external-target uncertainty. \\
Sensitivity &
Plausible target perturbations and departures from within-class exchangeability. \\
Reproducibility &
Software version or commit, random seeds, environment, frozen tables, and code that
recreates every reported number. \\
\bottomrule
\end{tabularx}
\end{table}

\section{Overall judgment}

\pkg\ is best described as \emph{prevalence-calibrated density-ratio weighting}.
That description is more precise than the expansion of the package name.  The
package does not infer unrestricted individual participation probabilities.  It fits
the smallest exponential modification of base weights that matches trusted
population moments, and it becomes true full-population IPW when the density-ratio
assumption in equation~\eqref{eq:identification-model} is correct.

For statistical genetics, the method is most persuasive as a generalization of
case--control ascertainment correction:
\begin{itemize}
\item it naturally combines a frame-based participation model with known register
prevalences;
\item it accommodates several diagnoses, joint moments, and stratum-specific margins;
and
\item it makes the correction inspectable through the tilt and weight diagnostics.
\end{itemize}
Its main risks are equally clear: marginal prevalence cannot restore an unidentified
joint distribution; exact anchors are tautologies; weak support creates extreme
weights; and current validation is synthetic.  The next scientific milestone is not
another zero calibration residual.  It is a prespecified external-cohort or held-out
register study showing improved quantities the fit was never told.

\clearpage
\begin{thebibliography}{99}

\bibitem{austin2009}
Austin, P. C. (2009).
Balance diagnostics for comparing the distribution of baseline covariates between
treatment groups in propensity-score matched samples.
\emph{Statistics in Medicine}, 28(25), 3083--3107.
\href{https://doi.org/10.1002/sim.3697}{doi:10.1002/sim.3697}.

\bibitem{bang2005}
Bang, H., and Robins, J. M. (2005).
Doubly robust estimation in missing data and causal inference models.
\emph{Biometrics}, 61(4), 962--973.
\href{https://doi.org/10.1111/j.1541-0420.2005.00377.x}
{doi:10.1111/j.1541-0420.2005.00377.x}.

\bibitem{chernozhukov2018}
Chernozhukov, V., Chetverikov, D., Demirer, M., Duflo, E., Hansen, C.,
Newey, W., and Robins, J. (2018).
Double/debiased machine learning for treatment and structural parameters.
\emph{The Econometrics Journal}, 21(1), C1--C68.
\href{https://doi.org/10.1111/ectj.12097}{doi:10.1111/ectj.12097}.

\bibitem{csiszar1975}
Csisz\'ar, I. (1975).
I-divergence geometry of probability distributions and minimization problems.
\emph{The Annals of Probability}, 3(1), 146--158.
\href{https://doi.org/10.1214/aop/1176996454}{doi:10.1214/aop/1176996454}.

\bibitem{dempster1950}
Dempster, E. R., and Lerner, I. M. (1950).
Heritability of threshold characters.
\emph{Genetics}, 35(2), 212--236.
\href{https://doi.org/10.1093/genetics/35.2.212}
{doi:10.1093/genetics/35.2.212}.

\bibitem{deville1992}
Deville, J.-C., and S\"arndal, C.-E. (1992).
Calibration estimators in survey sampling.
\emph{Journal of the American Statistical Association}, 87(418), 376--382.
\href{https://doi.org/10.1080/01621459.1992.10475217}
{doi:10.1080/01621459.1992.10475217}.

\bibitem{hainmueller2012}
Hainmueller, J. (2012).
Entropy balancing for causal effects: A multivariate reweighting method to produce
balanced samples in observational studies.
\emph{Political Analysis}, 20(1), 25--46.
\href{https://doi.org/10.1093/pan/mpr025}{doi:10.1093/pan/mpr025}.

\bibitem{hajek1971}
H\'ajek, J. (1971).
Comment on a paper by D. Basu.
In V. P. Godambe and D. A. Sprott (eds.),
\emph{Foundations of Statistical Inference}.
Holt, Rinehart and Winston.

\bibitem{hansen1982}
Hansen, L. P. (1982).
Large sample properties of generalized method of moments estimators.
\emph{Econometrica}, 50(4), 1029--1054.
\href{https://doi.org/10.2307/1912775}{doi:10.2307/1912775}.

\bibitem{horvitz1952}
Horvitz, D. G., and Thompson, D. J. (1952).
A generalization of sampling without replacement from a finite universe.
\emph{Journal of the American Statistical Association}, 47(260), 663--685.
\href{https://doi.org/10.1080/01621459.1952.10483446}
{doi:10.1080/01621459.1952.10483446}.

\bibitem{kott2010}
Kott, P. S., and Chang, T. (2010).
Using calibration weighting to adjust for nonignorable unit nonresponse.
\emph{Journal of the American Statistical Association}, 105(491), 1265--1275.
\href{https://doi.org/10.1198/jasa.2010.tm09016}
{doi:10.1198/jasa.2010.tm09016}.

\bibitem{kish1965}
Kish, L. (1965).
\emph{Survey Sampling}.
Wiley.

\bibitem{lee2011}
Lee, S. H., Wray, N. R., Goddard, M. E., and Visscher, P. M. (2011).
Estimating missing heritability for disease from genome-wide association studies.
\emph{The American Journal of Human Genetics}, 88(3), 294--305.
\href{https://doi.org/10.1016/j.ajhg.2011.02.002}
{doi:10.1016/j.ajhg.2011.02.002}.

\bibitem{lipton2018}
Lipton, Z. C., Wang, Y.-X., and Smola, A. J. (2018).
Detecting and correcting for label shift with black box predictors.
In \emph{Proceedings of the 35th International Conference on Machine Learning},
PMLR 80, 3122--3130.
\href{https://proceedings.mlr.press/v80/lipton18a.html}
{proceedings.mlr.press/v80/lipton18a.html}.

\bibitem{manski1977}
Manski, C. F., and Lerman, S. R. (1977).
The estimation of choice probabilities from choice based samples.
\emph{Econometrica}, 45(8), 1977--1988.
\href{https://doi.org/10.2307/1914121}{doi:10.2307/1914121}.

\bibitem{munafo2018}
Munaf\`o, M. R., Tilling, K., Taylor, A. E., Evans, D. M., and Davey Smith, G. (2018).
Collider scope: When selection bias can substantially influence observed
associations.
\emph{International Journal of Epidemiology}, 47(1), 226--235.
\href{https://doi.org/10.1093/ije/dyx206}{doi:10.1093/ije/dyx206}.

\bibitem{prentice1979}
Prentice, R. L., and Pyke, R. (1979).
Logistic disease incidence models and case-control studies.
\emph{Biometrika}, 66(3), 403--411.
\href{https://doi.org/10.1093/biomet/66.3.403}
{doi:10.1093/biomet/66.3.403}.

\bibitem{qin1998}
Qin, J. (1998).
Inferences for case-control and semiparametric two-sample density ratio models.
\emph{Biometrika}, 85(3), 619--630.
\href{https://doi.org/10.1093/biomet/85.3.619}
{doi:10.1093/biomet/85.3.619}.

\bibitem{qinlawless1994}
Qin, J., and Lawless, J. (1994).
Empirical likelihood and general estimating equations.
\emph{The Annals of Statistics}, 22(1), 300--325.
\href{https://doi.org/10.1214/aos/1176325370}
{doi:10.1214/aos/1176325370}.

\bibitem{robins1994}
Robins, J. M., Rotnitzky, A., and Zhao, L. P. (1994).
Estimation of regression coefficients when some regressors are not always observed.
\emph{Journal of the American Statistical Association}, 89(427), 846--866.
\href{https://doi.org/10.1080/01621459.1994.10476818}
{doi:10.1080/01621459.1994.10476818}.

\bibitem{saerens2002}
Saerens, M., Latinne, P., and Decaestecker, C. (2002).
Adjusting the outputs of a classifier to new a priori probabilities: A simple
procedure.
\emph{Neural Computation}, 14(1), 21--41.
\href{https://doi.org/10.1162/089976602753284446}
{doi:10.1162/089976602753284446}.

\bibitem{sargan1958}
Sargan, J. D. (1958).
The estimation of economic relationships using instrumental variables.
\emph{Econometrica}, 26(3), 393--415.
\href{https://doi.org/10.2307/1907619}{doi:10.2307/1907619}.

\bibitem{schoeler2023}
Schoeler, T., et al. (2023).
Participation bias in the UK Biobank distorts genetic associations and downstream
analyses.
\emph{Nature Human Behaviour}, 7, 1216--1227.
\href{https://doi.org/10.1038/s41562-023-01579-9}
{doi:10.1038/s41562-023-01579-9}.

\bibitem{vanalten2024}
van Alten, S., Domingue, B. W., Faul, J., Galama, T., and Marees, A. T. (2024).
Reweighting UK Biobank corrects for pervasive selection bias due to volunteering.
\emph{International Journal of Epidemiology}, 53(3), dyae054.
\href{https://doi.org/10.1093/ije/dyae054}{doi:10.1093/ije/dyae054}.

\bibitem{vanalten2025}
van Alten, S., Domingue, B. W., Faul, J., Galama, T., and Marees, A. T. (2025).
Correcting for volunteer bias in GWAS increases SNP effect sizes and heritability
estimates.
\emph{Nature Communications}, 16, article 3578.
\href{https://doi.org/10.1038/s41467-025-58684-8}
{doi:10.1038/s41467-025-58684-8}.

\bibitem{zhao2017}
Zhao, Q., and Percival, D. (2017).
Entropy balancing is doubly robust.
\emph{Journal of Causal Inference}, 5(1), article 20160010.
\href{https://doi.org/10.1515/jci-2016-0010}
{doi:10.1515/jci-2016-0010}.

\end{thebibliography}

\end{document}
